\documentclass[12pt]{article}

%--- Packages ------------------------------------------------------------------
\usepackage[margin=1in]{geometry}
\usepackage{amsmath,amssymb,amsthm}
\usepackage{setspace}
\usepackage{booktabs}
\usepackage[round,sort]{natbib}
\usepackage[hidelinks]{hyperref}

\onehalfspacing

%--- Theorem environments ------------------------------------------------------
\theoremstyle{plain}
\newtheorem{proposition}{Proposition}
\newtheorem{assumption}{Assumption}
\newtheorem{lemma}{Lemma}
\theoremstyle{remark}
\newtheorem{remark}{Remark}

%--- Title ---------------------------------------------------------------------
\title{Theoretical Framework\\
  \large Present Day, Present Time: AI Exposure and Labor Market Risk}
\author{}
\date{}

\begin{document}
\maketitle

% =============================================================================
\section{Theoretical Framework}
\label{sec:model}
% =============================================================================

\subsection{Setup}

Consider an economy in which production requires a continuum of tasks indexed
$i \in [0,1]$, ordered by their susceptibility to AI automation: tasks at the
low end are more easily performed by AI, while tasks at the high end require
judgment, communication, and contextual reasoning that AI currently struggles
to replicate.  This ordering maps naturally to the \citet{eloundou2024gpts}
beta scores: a high beta score indicates that a large share of an occupation's
tasks are susceptible to large language models, placing that occupation near
the low end of the $i$ spectrum.

Final output $Y$ aggregates tasks via a constant elasticity of substitution
(CES) technology:
\begin{equation}
  Y = \left[\int_0^1 y(i)^{\frac{\sigma-1}{\sigma}}\,di\right]^{\!\frac{\sigma}{\sigma-1}},
  \qquad \sigma > 0,
  \label{eq:ces}
\end{equation}
where $y(i)$ is the quantity of task $i$ used in production and $\sigma$ is
the elasticity of substitution across tasks.

Each task can be performed either by a human worker or by AI.  Human labor
performs any task at unit cost $w$ (the wage).  Performing task $i$ via AI
costs
\begin{equation}
  c^{\mathrm{AI}}(i) = \tau\cdot t(i),
  \label{eq:ai_cost}
\end{equation}
where $\tau > 0$ is the aggregate AI capability parameter---a lower $\tau$
reflects a more capable and cost-efficient AI system---and $t(i)$ is the
task-specific cost of automation, with $t'(i)>0$: harder tasks are more
expensive to automate.  We parameterize
\begin{equation}
  t(i) = e^{\alpha i}, \qquad \alpha > 0,
  \label{eq:cost_fn}
\end{equation}
so that automation costs rise smoothly and exponentially with task complexity.

\begin{assumption}[Partial Equilibrium]
  The wage $w$ is determined outside the model.  The economy is small and
  faces an infinitely elastic labor supply at wage $w$, following
  \citet{grossman2008trading}.  We abstract from general equilibrium wage
  determination.
\end{assumption}

% -----------------------------------------------------------------------------
\subsection{Task Assignment and the Automation Cutoff}
% -----------------------------------------------------------------------------

Since human labor and AI are perfect substitutes in the production of any
given task, in equilibrium each task is performed by whichever technology is
cheaper.  This yields a unique cutoff task $I^*$, defined by the indifference
condition:
\begin{equation}
  \tau\cdot t(I^*) = w.
  \label{eq:indiff}
\end{equation}
Tasks $i \in [0,I^*)$ satisfy $\tau t(i) < w$ and are performed by AI; tasks
$i \in [I^*,1]$ satisfy $\tau t(i) \geq w$ and are performed by human
workers.  Using \eqref{eq:cost_fn}, the cutoff is given in closed form by
\begin{equation}
  I^* = \frac{1}{\alpha}\ln\!\left(\frac{w}{\tau}\right).
  \label{eq:cutoff}
\end{equation}
The derivation is provided in Appendix~\ref{app:cutoff}.

A direct comparative static shows that AI improvement expands the set of
automated tasks:
\begin{equation}
  \frac{\partial I^*}{\partial\tau} = -\frac{1}{\alpha\tau} < 0.
  \label{eq:cs_tau}
\end{equation}
A fall in $\tau$---such as the deployment of GPT-4-class systems in
2023---raises $I^*$, pushing the automation frontier toward more complex
tasks.

% -----------------------------------------------------------------------------
\subsection{Displacement and Productivity Effects}
% -----------------------------------------------------------------------------

A rise in $I^*$ from $I^*_0$ to $I^*_1$ sets two opposing forces in motion,
as identified by \citet{grossman2008trading}.

\medskip
\noindent\textbf{Displacement effect.}  Workers performing tasks in the range
$[I^*_0, I^*_1]$ are directly displaced: their tasks are now cheaper to
automate.  The employment share of human labor falls by $\Delta I^* = I^*_1 -
I^*_0 > 0$.

\medskip
\noindent\textbf{Productivity effect.}  Automating low-$i$ tasks reduces the
cost of producing composite output $Y$, lowering its price index $P_Y$.  If
output demand is sufficiently elastic, lower $P_Y$ raises total demand for $Y$
and shifts out labor demand for the remaining human tasks $[I^*,1]$, partially
offsetting the displacement effect.

\medskip
The net effect on employment \emph{levels} is ambiguous in general and
constitutes an empirical question.  Section~\ref{sec:results} provides
evidence that the displacement effect dominates: higher AI exposure is
associated with elevated unemployment and lower wages, particularly following
the GPT-4 deployment shock of 2023.

% -----------------------------------------------------------------------------
\subsection{Education Heterogeneity}
\label{sec:model_het}
% -----------------------------------------------------------------------------

To generate testable predictions about differential impacts across education
groups, we introduce worker heterogeneity.  Consistent with the descriptive
evidence in \citet{eloundou2024gpts} and \citet{webb2020impact}, non-college
workers are disproportionately concentrated in tasks with high AI
substitutability (low $i$), while college-educated workers perform tasks
higher in the complexity spectrum.

\begin{assumption}[Education Sorting]
  Non-college workers supply tasks $i \in [0,\bar\imath]$ and college workers
  supply tasks $i \in [\bar\imath, 1]$, where $\bar\imath \in (I^*_0, 1)$ and
  $I^*_0$ is the pre-shock cutoff.
\end{assumption}

This two-type structure, while stylized, captures the core empirical
regularity.  It yields the following propositions, with proofs in
Appendix~\ref{app:proofs}.

\begin{proposition}[Differential Unemployment Effect]
  \label{prop:unemp}
  A fall in $\tau$ that raises $I^*$ into the interval $(I^*_0,\bar\imath)$
  displaces non-college workers while leaving college workers unaffected.  For
  any $\Delta\tau < 0$ such that $I^* < \bar\imath$, the displacement rate
  among non-college workers strictly exceeds that of college workers.
\end{proposition}

\begin{proposition}[Differential Wage Effect]
  \label{prop:wages}
  Workers performing tasks in the neighborhood of $I^*$ face the strongest
  downward wage pressure, as AI provides a close substitute at the task
  margin.  Since non-college workers are concentrated near $I^*$ while college
  workers are concentrated far from $I^*$, non-college wages experience larger
  compression following a fall in $\tau$.
\end{proposition}

\begin{remark}
  Propositions~\ref{prop:unemp} and~\ref{prop:wages} are consistent with AI
  capability in 2023--2025 being sufficient to automate many routine and
  semi-routine tasks (displacing non-college workers), but not yet the most
  complex professional tasks performed by college-educated workers.  The
  empirical results in Section~\ref{sec:results} support this timing: the
  interaction effects are statistically significant from 2023 onward and near
  zero in earlier years.
\end{remark}

% -----------------------------------------------------------------------------
\subsection{Task Bundling and Model Limitations}
\label{sec:bundling}
% -----------------------------------------------------------------------------

The framework above assumes \emph{atomistic} task assignment: workers freely
specialize in any subset of tasks from $[I^*,1]$ without constraint.  In
practice, occupations bundle tasks together, and workers must perform the
entire bundle rather than selecting individual tasks.  \citet{autor2025expertise}
show that this matters for the direction of wage effects: automation of
\emph{inexpert} tasks within a bundle raises the occupation's expertise
requirement and increases wages, while automation of \emph{expert} tasks
lowers the requirement and compresses wages.  Our model abstracts from this
dimension---the beta-score exposure measure captures the \emph{share} of tasks
exposed but not whether those tasks are the expert or inexpert ones within each
occupation's bundle.  We acknowledge this as a limitation, and regard the
expertise framework of \citeauthor{autor2025expertise} as a natural direction
for further work.

% =============================================================================
%  APPENDIX
% =============================================================================
\appendix

\section{Derivation of the Equilibrium Cutoff}
\label{app:cutoff}

Each task $i$ is performed by the lower-cost technology.  Human labor costs
$w$ per unit of any task.  AI performs task $i$ at cost $c^{\mathrm{AI}}(i) =
\tau e^{\alpha i}$.

\medskip
AI is cheaper whenever $\tau e^{\alpha i} < w$, i.e.,
\[
  e^{\alpha i} < \frac{w}{\tau}
  \implies
  i < \frac{1}{\alpha}\ln\!\left(\frac{w}{\tau}\right).
\]
Human labor is cheaper for $i > \frac{1}{\alpha}\ln(w/\tau)$.  The cutoff
task satisfies both conditions with equality:
\[
  I^* = \frac{1}{\alpha}\ln\!\left(\frac{w}{\tau}\right).
\]
For $I^* \in (0,1)$ we require $w/e^{\alpha} < \tau < w$, which we assume
holds at the pre-shock parameter values.

\medskip\noindent\textbf{Comparative static.}  Differentiating with respect to
$\tau$:
\[
  \frac{\partial I^*}{\partial\tau}
  = \frac{1}{\alpha}\cdot\frac{\partial}{\partial\tau}\bigl[\ln w - \ln\tau\bigr]
  = -\frac{1}{\alpha\tau} < 0.
\]
A fall in $\tau$ raises $I^*$: more tasks are automated as AI improves.
$\blacksquare$

\section{Proofs of Propositions}
\label{app:proofs}

\subsection*{Proof of Proposition~\ref{prop:unemp}}

Under Assumption 2, non-college workers perform tasks in $[I^*_0,\bar\imath]$
(tasks $[0,I^*_0]$ are already automated pre-shock).  After a fall in $\tau$,
the cutoff rises to $I^*_1 > I^*_0$.

Workers performing tasks in $[I^*_0, I^*_1]$ are displaced.  Provided
$I^*_1 \leq \bar\imath$ (the shock is moderate and does not reach the college
task range), all displaced workers are non-college.  The non-college
displacement rate is
\[
  d_{\mathrm{NC}} = \frac{I^*_1 - I^*_0}{\bar\imath - I^*_0} > 0,
\]
while the college displacement rate is $d_C = 0$.  Hence
$d_{\mathrm{NC}} > d_C$ for any $\Delta\tau < 0$.  $\blacksquare$

\subsection*{Proof of Proposition~\ref{prop:wages}}

In the neighborhood of $I^*$, AI provides a close substitute for human labor.
Competitive pressure from AI acts as an implicit price ceiling on the wages of
tasks near $I^*$: employers can replace workers in those tasks at cost
$\tau t(i) \approx w$ for $i$ just above $I^*$.  As $I^*$ rises, tasks
previously interior to the human task range now face direct AI competition,
compressing wages for workers in those tasks.

Under Assumption 2, non-college workers are concentrated in $[I^*_0,
\bar\imath]$---the lower end of the human task spectrum, immediately above
$I^*$.  College workers in $[\bar\imath, 1]$ are further from $I^*$ and
therefore face weaker competitive pressure.  It follows that non-college
workers experience (weakly) greater wage compression following a fall in
$\tau$.

\medskip
\noindent\textit{Note.}  Because wages are exogenous in the partial
equilibrium framework (Assumption 1), this proposition is best understood as a
directional prediction: given any positive relationship between AI competition
at the task margin and downward wage pressure, non-college workers bear a
disproportionate share.  The empirical specification in
Section~\ref{sec:empirical} tests this prediction directly.  $\blacksquare$

% =============================================================================
\bibliographystyle{plainnat}
\bibliography{references}

\end{document}
